\section{Methodology}
\label{ch:method}

\subsection{Overview}

Each month, a strategy decides a \emph{target} split from past data only; the engine then invests the instalment towards that target, rebalances if needed, charges costs and lets the month's market returns act on the holdings. Strategies are designed on 2010--2019 and tested on 2020--2025.

\subsection{The rolling (walk-forward) window}
\label{sec:window}

Strategies that use an optimiser need estimates of average returns, volatilities and correlations.
These are always computed from the \textbf{previous 120 months} (10 years). To decide month $t$, the
optimiser sees months $t-120$ to $t-1$ only. When the window moves forward by a month, the newest
month joins and the oldest drops out, so the window always holds 120 months and never contains
future data.

Because the data starts in February 2000, the first decision that has a full window is for
February 2010 (window Feb~2000 -- Jan~2010). All strategies, including those that do not need a
window, are therefore run from February 2010 so that they are compared over the same months.
Strategies 4--7 re-fit their target once a year, every February; the Personal SIP re-fits every
month.

\subsection{The SIP engine}
\label{sec:engine}

Each strategy is defined by three things: how its target is set, how the instalment is split, and
when it rebalances. The engine runs the same five steps every month for every strategy:

\begin{enumerate}
  \item \textbf{Instalment.} \Rs10{,}000 is invested on the first of the month.
  \item \textbf{Split the instalment.} Either \emph{pro rata} ($w_i \times$ amount) or by the
        \emph{smart split} described below.
  \item \textbf{Rebalance if the rule says so.} Either never, once a year in January
        (\emph{calendar}), or whenever any asset is more than 5 percentage points away from its
        target (\emph{band}). Rebalancing trades every holding back to the target.
  \item \textbf{Costs.} 0.1\% of the value traded is deducted.
  \item \textbf{Market moves.} Each holding grows by its asset's return for the month.
\end{enumerate}

\paragraph{Smart split.} New money is sent first to the assets that are below their target. Let
$h_i$ be the current holdings, $A$ the instalment, $T = \sum_i h_i + A$ the total after investing,
and $g_i = \max(w_i T - h_i, 0)$ the gap of asset $i$, with $G = \sum_i g_i$. Then the purchase of
asset $i$ is
\begin{equation}
  b_i =
  \begin{cases}
    g_i \,/\, G \times A & \text{if } G \ge A \quad\text{(share the money in proportion to the gaps)},\\[2pt]
    g_i + w_i\,(A - G) & \text{if } G < A \quad\text{(fill every gap, split the rest by target)}.
  \end{cases}
  \label{eq:smart}
\end{equation}
Since $\sum_i (w_i T - h_i) = T - \sum_i h_i = A$, the positive gaps always add up to at least the
instalment, so $G \ge A$ and the first case always applies; the second case is kept only as a
numerical safety net. If no asset is above its target, $G = A$ and every gap is filled exactly.

\paragraph{Band rebalancing.} After the purchase, the drift is
$d = \max_i \left| \frac{h_i + b_i}{T} - w_i \right|$. If $d > 5\%$ every holding is traded to
$w_i T$. A band rather than a calendar rule avoids trading for small deviations, as recommended by
\citet{jaconetti2010}.

\subsection{Performance measures}
\label{sec:metrics}

\begin{itemize}
  \item \textbf{Time-weighted return (TWR).} The strategy's return in month $t$, measured after the
        instalment and before the next one, so that it reflects the strategy and not the timing of
        cash flows.
  \item \textbf{XIRR.} The single annual rate $x$ that makes the value of all instalments, grown at
        $x$, equal to the final portfolio value. It answers ``what yearly return did my SIP money
        earn?'' and is found numerically.
  \item \textbf{Volatility.} $\sigma = \text{std}(\text{TWR}) \times \sqrt{12}$. Multiplying by
        $\sqrt{12}$ converts a monthly standard deviation to a yearly one, since variances of
        independent months add up.
  \item \textbf{Sharpe ratio versus liquid.}
        $S = 12 \times \text{mean}(\text{TWR}_t - r_{\text{Liquid},t}) \,/\, \sigma$: return above
        the investor's real safe alternative, per unit of risk \citep{sharpe1994}.
  \item \textbf{Maximum drawdown.} The worst fall of the growth-of-\Rs1 index from a previous peak.
\end{itemize}

\subsection{The utility score $U$}
\label{sec:utility}

To compare strategies with one number that balances profit against risk for a given investor, the
mean-variance utility is used \citep{markowitz1952,levy1979,bkm2021}:
\begin{equation}
  U = \underbrace{12 \times \text{mean}(\text{TWR}_t - r_{\text{Liquid},t})}_{\text{return above liquid}}
      \; - \; \frac{\lambda}{2}\,\underbrace{\left(\text{std}(\text{TWR}) \sqrt{12}\right)^2}_{\text{variance}} .
  \label{eq:U}
\end{equation}

\begin{itemize}
  \item The first term is the reward: how much more the strategy earned than the liquid fund.
  \item The second term is a penalty for variability. Volatility is squared, so doubling the swings
        quadruples the penalty, in line with the observation that large losses hurt far more than
        small ones.
  \item $\lambda$ is the investor's risk aversion \citep{pratt1964}. Three profiles were fixed in
        advance: \textbf{aggressive} $\lambda=2$, \textbf{moderate} $\lambda=4$ and
        \textbf{conservative} $\lambda=8$.
  \item $U$ is a certainty equivalent: $U = 1.36\%$ means the investor is as satisfied with the
        risky strategy as with a guaranteed 1.36\% a year above the liquid fund.
\end{itemize}

\paragraph{Connection to Lagrange multipliers.} Maximising $U$ is Markowitz's problem of maximising expected return for a given variance; $\lambda$ is the Lagrange multiplier on the variance constraint, i.e.\ the price the investor puts on risk.

\paragraph{Example.} The 60/20/20 SIP on 2010--2019 earned 2.91\% a year above liquid with 8.81\%
volatility. For a moderate investor, $U = 2.91\% - 2 \times (8.81\%)^2 = 2.91\% - 1.55\% = 1.36\%$.

\subsection{The strategies}

All strategies hold the same three assets and run in the same engine. Table~\ref{tab:strategies}
summarises them.

\begin{table}[H]
\centering
\caption{The strategies compared}
\label{tab:strategies}
\small
\begin{tabularx}{\textwidth}{@{}clXll@{}}
\toprule
\# & Name & Target & Instalment & Rebalance \\
\midrule
1 & Equity SIP & 100\% Nifty & pro rata & never \\
2 & Equal-weight SIP & $\tfrac13$ each & pro rata & never \\
3 & 60/20/20 SIP & 60\% Nifty, 20\% gold, 20\% liquid & pro rata & every January \\
4 & Risk-parity SIP & equal risk contributions, refit every Feb & smart & 5\% band \\
5 & Max-Sharpe SIP & maximum Sharpe ratio (vs 0\%), refit every Feb & smart & 5\% band \\
6 & Optimized SIP & $\tfrac12$ (4) + $\tfrac12$ (5) & smart & 5\% band \\
6b & Fixed Optimized SIP & liquid fixed at 10\%; Nifty/gold by $\tfrac12$ RP + $\tfrac12$ max Sharpe vs liquid & smart & 5\% band \\
7 & Max-Utility SIP & maximises $U$ for the investor's $\lambda$, refit every Feb & smart & 5\% band \\
8 & Personal SIP & liquid cushion by $\lambda$; rest as 6b; refit monthly & smart & 5\% band \\
\bottomrule
\end{tabularx}
\end{table}

In Strategies 4--7 every asset is kept between 10\% and 70\% of the portfolio, so no optimiser can
put everything in one asset.

\subsubsection{Risk parity (Strategy 4)}
\label{sec:rp}

The \emph{risk contribution} of asset $i$ is its weight times its marginal contribution to variance
\citep{maillard2010,qian2006}:
\begin{equation}
  RC_i = w_i\,(\Sigma w)_i, \qquad \sum_i RC_i = w^\top\Sigma w .
  \label{eq:rc}
\end{equation}
For two assets with variances $V_N, V_G$ and covariance $C$, this splits the portfolio variance
$w_N^2V_N + 2w_Nw_GC + w_G^2V_G$ into each asset's own term plus half of the shared term:
\begin{equation}
  RC_N = w_N\left(w_N V_N + w_G C\right), \qquad RC_G = w_G\left(w_G V_G + w_N C\right).
\end{equation}
Risk parity chooses weights so that all contributions are equal. With three assets this is solved
numerically by minimising $\sum_i (RC_i - \overline{RC})^2$ (sequential least squares, SLSQP). With
two assets, setting $RC_N = RC_G$ cancels the shared terms and gives a closed form:
\begin{equation}
  w_N^2 V_N = w_G^2 V_G \;\Rightarrow\; w_N\sigma_N = w_G\sigma_G \;\Rightarrow\;
  w_N = \frac{\sigma_G}{\sigma_N + \sigma_G}.
  \label{eq:rp2}
\end{equation}
The more volatile asset receives the smaller weight. Risk parity uses no expected returns.

\subsubsection{Maximum Sharpe ratio (Strategy 5)}
\label{sec:maxsharpe}

The maximum Sharpe portfolio maximises expected return per unit of volatility:
\begin{equation}
  \max_w \; \frac{w^\top \mu}{\sqrt{w^\top \Sigma w}} \quad \text{subject to } \textstyle\sum_i w_i = 1,\; 0.1 \le w_i \le 0.7 .
  \label{eq:ms}
\end{equation}
For two assets with $w$ in Nifty, the numerator is a simple weighted average,
$m = w\mu_N + (1-w)\mu_G$, while the variance
$v = w^2V_N + 2w(1-w)C + (1-w)^2V_G$ is not. When the two assets are weakly or negatively correlated
the middle term is small or negative, so a mix can have lower volatility than either asset alone,
and therefore a higher Sharpe ratio. Strategy 5 uses raw returns, i.e.\ excess over 0\%.

\subsubsection{The Optimized SIP (Strategy 6)}

The average of the risk-parity and maximum-Sharpe weights. The two components fail in different
ways: maximum Sharpe depends on noisy average returns \citep{chopra1993}, while risk parity ignores
returns altogether. Averaging them is a simple way to hedge between the two.

\subsubsection{The Fixed Optimized SIP (Strategy 6b)}
\label{sec:fixed}

Section~\ref{sec:res-2010} shows that Strategies 4--6 put 70\%, the maximum allowed, into the liquid
fund. Its volatility is about 0.5\%, so measured against a 0\% rate it looks nearly risk-free with
a large return. The correction treats liquid as cash, following \citet{tobin1958} and
\citet{sharpe1994}:
\begin{itemize}
  \item The liquid share is fixed at $L$.
  \item The rest is split between Nifty and gold by $\tfrac12$ risk parity $+\ \tfrac12$ maximum
        Sharpe, where maximum Sharpe uses returns \emph{above liquid},
        $\mu_i - \mu_{\text{Liquid}}$.
  \item Nifty and gold each stay between 10\% and 70\% of the \emph{whole} portfolio. Inside the
        risky part, which is $1-L$ of the money, this means a Nifty share between
        $\ell = 0.1/(1-L)$ and $\min\!\left(0.7/(1-L),\, 1 - \ell\right)$.
\end{itemize}
$L$ was chosen from 10\%, 20\% and 30\% by the highest Sharpe ratio versus liquid on 2010--2019 only:
10\% won (0.291 against 0.278 and 0.276) and was frozen before any later data was examined.

\subsubsection{The Max-Utility SIP (Strategy 7)}

Strategy 7 maximises the utility of Eq.~\eqref{eq:U} directly, using the window's yearly estimates:
\begin{equation}
  \max_w \; w^\top(\mu - \mu_{\text{Liquid}}) - \frac{\lambda}{2}\, w^\top \Sigma w ,
  \qquad \textstyle\sum_i w_i = 1,\; 0.1 \le w_i \le 0.7 .
\end{equation}
Because returns are measured above liquid, the liquid share is decided by $\lambda$ itself: on the
February 2010 window, $\lambda = 1, 2, 4$ held the minimum 10\% in liquid, while $\lambda = 8$ held
47\% and $\lambda = 16$ the maximum 70\%.

\subsubsection{The Personal SIP (Strategy 8)}
\label{sec:personal}

The Fixed Optimized SIP gives every investor the same 10\% liquid share, so the investor's risk
aversion does not change its advice. The Personal SIP applies Tobin's separation directly:
\begin{enumerate}
  \item \textbf{How much to keep safe} is personal: a liquid cushion $L$ set by $\lambda$.
  \item \textbf{How to split the rest} is common: Nifty and gold as in Strategy 6b, re-fitted every
        month on the latest 120 months.
\end{enumerate}
The target is therefore
\begin{equation}
  w = \Big[(1-L)\,n,\;\; (1-L)(1-n),\;\; L\Big], \qquad
  n = \tfrac12\, n_{\text{RP}} + \tfrac12\, n_{\text{MS}} ,
  \label{eq:personal}
\end{equation}
where $n_{\text{RP}}$ is the Nifty share of the risky part from Eq.~\eqref{eq:rp2} and $n_{\text{MS}}$
the Nifty share that maximises Eq.~\eqref{eq:ms} on returns above liquid, both clipped to the bounds
of Section~\ref{sec:fixed}.

\paragraph{Choosing the cushion.} For each profile, cushions of 10\%, 20\%, \dots, 70\% were each run
through the full engine on Feb~2010 -- Dec~2019, scored with $U$ for that profile's $\lambda$, and
the cushion with the highest $U$ was frozen (Section~\ref{sec:res-personal}). $U$ is therefore used
once, to choose the cushion; it is not recomputed when the rule is applied each month.

\subsection{Experimental design}
\label{sec:design}

To avoid a rule that only looks good on the data it was built on \citep{bailey2014}, the data is
split in time. Feb~2000 -- Jan~2010 serves only as the first look-back window. All strategies are
compared, and the liquid share of 6b and the cushions of the Personal SIP are chosen, on
Feb~2010 -- Dec~2019. Every rule is then run unchanged on Jan~2020 -- Dec~2025, with the rule for
picking the winner fixed in advance. The test-period ranking is also rebuilt from alternative funds,
and the simple strategies are recomputed with independent code.
